\documentclass[aspectratio=169]{beamer}

\usetheme{Madrid}
\usecolortheme{default}

\usepackage[utf8]{inputenc}
\usepackage[spanish]{babel}
\usepackage{booktabs}
\usepackage{amsmath,amssymb}
\usepackage{tikz}
\usetikzlibrary{arrows.meta,positioning}

\title[SigmaBurger]{SigmaBurger}
\subtitle{Prompt Injection en Chatbots de Atención al Público}
\author[Gonzalez, Salinas]{Felipe Gonzalez \and Pedro Salinas}
\institute[IIC2531]{IIC2531 --- Proyecto Final}
\date{Noviembre 2026}

\setbeamertemplate{navigation symbols}{}
\setbeamertemplate{caption}[numbered]

\AtBeginSection[]{
  \begin{frame}{Contenidos}
    \tableofcontents[currentsection]
  \end{frame}
}

\begin{document}

\begin{frame}
  \titlepage
\end{frame}

\begin{frame}{Contenidos}
  \tableofcontents
\end{frame}

\section{Introducción}

\begin{frame}{Un incidente real}
  \begin{quote}
    ``Vendimos un vehículo nuevo 2024 por \$1 dólar.''
  \end{quote}
  \begin{flushright}
    --- chatbot de atención al cliente de un concesionario Chevrolet,\\
    diciembre 2023, engañado con un \emph{prompt injection}.
  \end{flushright}
  \vfill
  \begin{block}{Pregunta de investigación}
    ¿Qué defensas reducen el éxito de ataques de inyección directa en un
    chat bot, y a qué costo?
  \end{block}
\end{frame}

\begin{frame}{El problema: un solo canal para todo}
  \centering
  \vspace{2mm}
  \begin{tikzpicture}[
      x=1cm,
      y=1cm,
      card/.style={draw, rounded corners=3pt, align=center, text width=3cm,
        minimum height=1.3cm, line width=0.9pt, inner sep=8pt},
      flow/.style={-{Stealth[length=2.5mm]}, line width=1pt, draw=gray!75},
      every node/.style={font=\small}
    ]
    \node[card, fill=blue!8, draw=blue!55] (system) at (0,1.35)
      {\textbf{System prompt}\\\footnotesize Reglas del negocio};
    \node[card, fill=orange!8, draw=orange!75!black, dashed] (user) at (0,-1.35)
      {\textbf{Mensaje del usuario}\\\footnotesize Entrada no confiable};
    \node[card, text width=3.2cm, minimum height=2cm, fill=violet!9, draw=violet!65]
      (context) at (4.7,0)
      {\textbf{Contexto compartido}\\[2pt]\Large LLM};
    \node[card, text width=3cm, minimum height=2cm, fill=teal!8, draw=teal!60!black]
      (tools) at (9.2,0)
      {\textbf{Herramientas}\\[2pt]\footnotesize Menú · pedidos · descuentos\\[3pt]
       \scriptsize Sin validación de argumentos por defecto};

    \draw[flow] (system.east) to[out=0,in=155] (context.155);
    \draw[flow, draw=orange!80!black] (user.east) to[out=0,in=205] (context.205);
    \draw[flow] (context.east) -- node[above=4pt, font=\scriptsize, text=gray!80!black]
      {argumentos} (tools.west);
  \end{tikzpicture}
  \vspace{3mm}
  \vfill
  {\footnotesize\color{gray!80!black}Referencia: OWASP LLM01:2025 --- Prompt Injection}
\end{frame}

\begin{frame}{SigmaBurger: el agente bajo prueba}
  \begin{itemize}
    \item Agente de pedidos de comida rápida.
    \item Implementado en Python, sobre la API de \emph{Chat Completions}
      de OpenAI con \emph{function calling}.
    \item Herramientas expuestas al modelo:
      \begin{itemize}
        \item \texttt{get\_menu}
        \item \texttt{create\_order}
        \item \texttt{apply\_discount}
      \end{itemize}
    \item El atacante interactúa como un cliente más, por el canal normal
      de conversación (inyección directa).
  \end{itemize}
\end{frame}

\begin{frame}{Qué priorizamos}
  \begin{table}
    \begin{tabular}{@{}ll@{}}
      \toprule
      Propiedad & Alcance \\
      \midrule
      \textbf{Integridad} & En alcance: pedidos, precios, descuentos \\
      Confidencialidad & Fuera de alcance \\
      Disponibilidad & Fuera de alcance \\
      \bottomrule
    \end{tabular}
  \end{table}
  \vfill
\end{frame}

\section{Metodología}

\begin{frame}{Qué cuenta como violación}
  \begin{block}{Éxito efectivo}
    Queda un pedido con más de 10\% de descuento, o con un precio distinto
    al del menú. \textbf{Daño real.}
  \end{block}
  \begin{block}{Intento}
    El modelo invoca la herramienta con argumentos indebidos, pero el
    código lo bloquea antes de ejecutarse. \textbf{Sin daño.}
  \end{block}
\end{frame}

\begin{frame}{Catálogo de ataques: 24 ataques, 4 familias}
  \begin{table}
    \begin{tabular}{@{}lc@{}}
      \toprule
      Familia & N.\textordmasculine{} de ataques \\
      \midrule
      Suplantación & 7 \\
      Persuasión & 9 \\
      Evasión & 6 \\
      Directo & 2 \\
      \midrule
      Total & 24 \\
      \bottomrule
    \end{tabular}
  \end{table}
\end{frame}

\begin{frame}{Defensas evaluadas: D0 a D7}
  \begin{table}
    \small
    \begin{tabular}{@{}ll@{}}
      \toprule
      Código & Descripción \\
      \midrule
      D0 & Sin defensa \\
      D1 & Prompt reforzado \\
      D2 & Spotlighting \\
      D3 & Clasificador de entrada \\
      D5 & Validación en código \\
      D6 & D1 + D3 + D5 \\
      D7 & D1 + D5 \\
      \bottomrule
    \end{tabular}
  \end{table}
\end{frame}

\begin{frame}{Diseño experimental}
  \begin{itemize}
    \item \textbf{2 modelos}: \texttt{gpt-4o-mini} y \texttt{gpt-6-luna}.
    \item \textbf{Ataques $\times$ repeticiones}: $24 \times 10$ para
      gpt-4o-mini, $24 \times 5$ para gpt-6-luna.
    \item \textbf{12 conversaciones legítimas} por defensa, para medir
      falsos positivos.
    \item \textbf{Costo}: tokens consumidos por conversación legítima,
      normalizado contra D0.
  \end{itemize}
\end{frame}

\section{Resultados}

\begin{frame}{De 30\% a casi 0\%: tasa de éxito efectivo (ASR)}
  \begin{table}
    \small
    \centering
    \begin{tabular}{@{}lcccccccc@{}}
      \toprule
      Modelo & D0 & D1 & D2 & D3 & D5 & D6 & D7 \\
      \midrule
      gpt-4o-mini & 30.4\% & 0\% & 5.4\% & 0\% & 0\% & 0\% & 0\% \\
      gpt-6-luna & 35.8\% & 0\% & 20.0\% & 0\% & 0\% & 0\% & 0\% \\
      \bottomrule
    \end{tabular}
  \end{table}
  \vfill
  \begin{alertblock}{Hallazgo principal}
    D5 es la única defensa que, sola, lleva el daño real a 0\%
    (IC 95\% por bootstrap).
  \end{alertblock}
\end{frame}

\begin{frame}{Huella de vulnerabilidad por familia (D0)}
  \begin{table}
    \small
    \centering
    \begin{tabular}{@{}lcccc@{}}
      \toprule
      Modelo & Directo & Suplant. & Persuas. & Evasión \\
      \midrule
      gpt-4o-mini & 0\% & 64.3\% & 22.2\% & 13.3\% \\
      gpt-6-luna & 10.0\% & 22.9\% & 48.9\% & 40.0\% \\
      \bottomrule
    \end{tabular}
  \end{table}
  \vfill
  \begin{itemize}
    \item \texttt{gpt-4o-mini} falla sobre todo a suplantación.
    \item \texttt{gpt-6-luna} falla sobre todo a persuasión y evasión.
  \end{itemize}
\end{frame}

\begin{frame}{El costo de defenderse (gpt-4o-mini)}
  \begin{table}
    \small
    \centering
    \begin{tabular}{@{}lccc@{}}
      \toprule
      Defensa & Daño real & Falsos positivos & Costo (tokens, $\times$ D0) \\
      \midrule
      D0 & 30.4\% & 0\% & 1.0$\times$ \\
      D1 & 0\% & 0\% & 2.35$\times$ \\
      D2 & 5.4\% & 0\% & 1.21$\times$ \\
      D3 & 0\% & 30.0\% & 0.98$\times$ \\
      D5 & 0\% & 0\% & 1.0$\times$ \\
      D6 & 0\% & 30.8\% & 2.05$\times$ \\
      D7 & 0\% & 0\% & 2.4$\times$ \\
      \bottomrule
    \end{tabular}
  \end{table}
  \vfill
\end{frame}

\begin{frame}{El costo oculto de un clasificador}
  \begin{itemize}
    \item Solo D3 y D6, que incluyen un clasificador de entrada, rechazan
      a clientes legítimos.
    \item gpt-4o-mini: 30.0\% (D3) y 30.8\% (D6) de falsos positivos.
    \item gpt-6-luna: 3.3\% (D3), 0\% en el resto.
  \end{itemize}
\end{frame}

\begin{frame}{Hallazgo inesperado}
  \begin{alertblock}{Sin ningún ataque de por medio}
    El modelo creó un segundo pedido y le aplicó 50\% de descuento que
    nadie solicitó.
  \end{alertblock}
  \begin{itemize}
    \item 3 casos observados.
    \item 0 ataques involucrados.
    \item Solo D5 (validación en código) lo detiene.
  \end{itemize}
\end{frame}



\begin{frame}{Conclusión}
  \begin{alertblock}{}
    \centering
    \Large
    La única defensa que elimina el daño real es la que no depende de que
    el modelo se porte bien.
  \end{alertblock}
  \vfill
  \textbf{Trabajo futuro:}
  \begin{itemize}
    \item Ataques de confidencialidad.
    \item Clasificador con otro modelo.
    \item Atacante adaptativo.
    \item Más repeticiones donde los intervalos de confianza se superponen.
  \end{itemize}
\end{frame}

\section{Cierre}

\begin{frame}
  \centering
  \Large Gracias. ¿Preguntas?
  \vfill
  \normalsize \url{github.com/Pedr0sit0s/proyecto_seguridad}
\end{frame}

\end{document}
